\chapter{Scope, status, and inverse conventions}
\label{chap:inverse-scope-status}

\section{What this volume consolidates}
\label{sec:scope}

The source corpus consisted of six independently delivered packages:
\emph{Inverse \(q\)-Analog Functions},
\emph{Inverse \(q\)-Analog Jet Atlas},
\emph{Inverse \(q\)-Analogs: Extended Report},
two reports entitled \emph{Inverse \(q\)-Analogs}, and
\emph{\(q\)-Pochhammer and \(q\)-Binomial Expansions}.  They treated largely
the same inverse maps in different orders, repeatedly reproved the same
reversion formulas, and occasionally assigned incompatible hypotheses or
coefficients to the same claim.  The present volume is an editorial and
mathematical synthesis, not a mechanical concatenation.

The forward-expansion report supplies the strongest systematic engine for
finite products, Gaussian coefficients, roots of unity, and the
\(q\to1\) regimes.  The jet atlas supplies the most useful all-order inverse,
Bell-polynomial, mixed-parameter, and Lagrange--Good formulations.  The
function report supplies the clearest branchwise narrative.  The extended and
two earlier inverse reports contribute certification arguments, special
regimes, and the Fabius--Rvachev applications.  Repeated material appears once,
with the strongest correct hypotheses and the most transparent proof.

\section{Status vocabulary}
\label{sec:status-vocabulary}

Three statuses are used and are not interchangeable.

\begin{description}
\item[Lean-proved.]
The statement has a named, compiler-validated declaration in
\repo.  Its exact module and declaration name are recorded in the reference
appendix.
\item[Human-proved frontier result.]
The statement is proved completely in this volume but does not yet have an
exact Lean counterpart.  It belongs in this research-frontier document rather
than in the primary formalization-backed exposition.
\item[Conjecture.]
The statement is plausible and mathematically meaningful, but the argument
currently has a genuine missing step.  Conjectures are isolated in explicitly
named \emph{Conjecture} environments.  A numerical experiment, symbolic check,
or manuscript label is never treated as a proof.
\end{description}

Definitions, notation, and descriptions of an algorithm are not themselves
assertions of correctness.  Whenever correctness, convergence, error, or
complexity is claimed, it is stated as a theorem or proposition and proved.

\section{Basic normalizations}
\label{sec:normalizations}

For \(n\in\NaturalNumbers\), complex parameters \(a,q\), and an empty product equal to one,
the finite shifted factorial is
\[
  \QPochhammer{a}{q}{n}:=\prod_{j=0}^{n-1}(1-aq^j).
\]
When \(|q|<1\), the infinite shifted factorial is
\[
  \QPochhammer{a}{q}{\infty}:=\prod_{j=0}^{\infty}(1-aq^j).
\]
The convergence and holomorphy assertions implicit in the latter notation are
proved in \cref{sec:infinite-normal-convergence}; outside \(|q|<1\), the symbol
is used only after an explicit continuation or reciprocal normalization has
been specified.

For integers \(0\le k\le n\), the Gaussian coefficient is normalized by
\[
  \GaussianBinomial{n}{k}{q}
  :=\frac{\QPochhammer{q}{q}{n}}
  {\QPochhammer{q}{q}{k}\QPochhammer{q}{q}{n-k}},
\]
with its polynomial continuation at points where this quotient is formally
indeterminate.  This convention has degree \(k(n-k)\), constant coefficient
one, and satisfies the reciprocity law
\[
  \GaussianBinomial{n}{k}{q}
  =q^{k(n-k)}\GaussianBinomial{n}{k}{q^{-1}}
\]
for \(q\ne0\).  A proof independent of quotient cancellation appears in
\cref{sec:gaussian-reciprocity}.

We use
\[
  q=\EulerE^{-t}\qquad(t\to0)
\]
for logarithmic expansions near \(q=1\), and \(h=q-1\) only when an ordinary
Taylor coordinate is genuinely preferable.  Conversion between the two is
performed by formal composition, never by silently replacing \(t\) with
\(-h\).  At a root of unity \(\zeta\), radial limits are written
\(q=\zeta\EulerE^{-t}\) with \(t\downarrow0\); this specifies a sectorial or radial
asymptotic branch and does not assert the existence of an analytic germ across
the unit circle.

\section{What an inverse symbol means}
\label{sec:inverse-data}

An expression such as \(a=a(y;q)\) or \(q=q(y;a)\) is shorthand, not a
globally defined function.  Every inverse used below records the following
data:

\begin{enumerate}
\item the forward map and which parameters are held fixed;
\item the parameter being recovered;
\item the base point, limiting regime, or real interval;
\item the local sheet or global monotone branch;
\item the logarithm or fractional-power determination when one is needed.
\end{enumerate}

Thus a scalar equation with several free parameters defines a level set, not a
simultaneous inverse.  Several parameters can be recovered only after enough
independent observables have been supplied and the corresponding Jacobian has
been shown invertible.  This elementary distinction removes a substantial
fraction of the apparent disagreements among the source reports.

\section{Corrections made during consolidation}
\label{sec:corrections}

The synthesis corrects, rather than propagates, the following recurrent
defects from the source corpus:

\begin{itemize}
\item for \(0<q<1\), the real zero lattice of
  \(a\mapsto\QPochhammer{a}{q}{\infty}\) is ordered
  \(1=q^0<q^{-1}<q^{-2}<\cdots\);
\item a centered-palindromic inverse theorem requires nonzero equal endpoint
  coefficients, not palindromy alone;
\item continuous lower-parameter Gaussian inversion must stay on the branch
  where the target lies strictly between the endpoint value and the central
  maximum;
\item zeros introduced by a cyclotomic factor are simple unless a separate
  factor collision is proved; multiplicity cannot be inferred from the
  factor's mere presence;
\item analytic-germ language at a natural boundary is replaced by radial or
  sectorial asymptotics with an explicit approach region;
\item Newton and residual certificates include the differentiability,
  simple-root, derivative-separation, and domain-control hypotheses on which
  their conclusions depend;
\item multivariate Lagrange--Good product formulas are used only for systems
  in the required product form; arbitrary locally invertible systems instead
  use the general total-degree recursion;
\item conjectures contradicted by elementary counterexamples are removed
  rather than relabelled.
\end{itemize}

Later chapters give the corrected formulas and proofs at the point of use.
